% displayname : Équations de Maxwell
\documentclass[12pt,fancy]{cours}
\Chapit{5}{5}
\begin{document}
\maketitre{
Nous abordons maintenant les équations complètes de Maxwell, en régime dépendant du temps, où l'évolution des champs électrique et magnétique est couplée.
}{}


%%%%%%%%%%%%%%%%%%%%%%%%%%%%%%%%%%%%%%%%%%
\section{Postulats de l’électromagnétisme}

\subsection{Charges et courants électriques}

\parag{Le champ électromagnétique} Nous avons vu qu'on pouvait interpréter les phénomènes électrostatiques et magnétostatiques  en introduisant deux champs indépendants : le champ électrostatique $\Vec{E}(M)$, produit par les charges électriques ; et le champ magnétostatique $\Vec{B}(M)$, produit par les courants électriques. Les phénomènes d'induction que vous avez vus en $1^{\rm ere}$ année sont venus troubler ce découplage apparent entre magnétisme et électricité. En régime variable, nous devons décrire les phénomènes magnétiques et électriques par un champ indissociable :
% \begin{center}
  \eqbox{
    \text{ le champ électromagnétique } \left(\Vec{E}(M,t),\Vec{B}(M,t)\right).
    }
% \end{center}
\begin{remarque}
Les champs électrique $\Vec{E}_{\Rcal}(M,t)$ et magnétique $\Vec{B}_{\Rcal}(M,t)$ dépendent en toute rigueur du référentiel $\mc{R}$, même si en pratique on ne précisera pas le référentiel d’étude.
\end{remarque}

\parag{Action d’un champ électromagnétique sur une charge}
Un champ électrostatique $\Vec{E}(M)$ exerce sur une charge ponctuelle $q$ au point $M$ une force
\begin{equation}
\Vec{F}_{\rm e} = q \Vec{E}(M).
\end{equation}
Cette contribution électrique garde la même expression pour un particule immobile dans un champ électrique variable, d’expression $\Vec{F}_{\rm e} = q \Vec{E}(M,t)$.
Considérons maintenant cette même charge ponctuelle placée au point $M$ à l'instant $t$, qui se déplace à la vitesse $\vR(t)$ dans le référentiel $\mc{R}$, dans lequel règne un champ électromagnétique. Pour ne pas alourdir les notations, dans toute la suite nous omettrons la mention du référentiel $\mc{R}$ et du point d'espace-temps $(M,t)$. Il s’exerce sur cette charge une force appelée \imp{force de Lorentz}, d’expression
\begin{theorem}[Postulat : force de Lorentz]\label{form:}
  Une charge ponctuelle $q$ de vitesse $\vv{v}$ placée dans un champ électromagnétique $(\vv{E},\vv{B})$ subit la force
\begin{equation*}
\vv{F}=q\left(\vv{E}+\vv{v}\pv\vv{B}\right).
\end{equation*}
\end{theorem}
\noindent
La force de Lorentz comprend donc deux contributions : 
\begin{liste}
  \item une force électrique $\Vec{F}_{\rm e} = q\Vec{E}$ ;
\item une force magnétique $\Vec{F}_{\rm m} = q\Vec{v} \pv \Vec{B}$. 
\end{liste}
\begin{remarque}
  Cela implique qu’un champ électrique (en $\SI{}{V.m^{-1}}$) est homogène à une vitesse fois un champ magnétique, soit $\SI{1}{V.m^{-1}} = \SI{1}{m.s^{-1}.T}$, soit $\SI{1}{V} = \SI{1}{m^2.T.s^{-1}}$. Une tension est donc homogène à un flux magnétique par unité de temps, en accord avec la loi de l’induction.
\end{remarque}

\begin{exemple}{Force de Laplace}
  On considère une assemblée de $n$ particules par unité de volume, de charge $q$ et de vitesse $\Vec{v}$, plongée dans un champ électromagnétique $(\Vec{E},\Vec{B})$. 
  \begin{quest}
    \item	Déterminer la résultante de la force de Lorentz qui s’exerce sur l’ensemble des particules d’une volume élémentaire $\d V$. En déduire la force de Lorentz volumique.
    \item Quelle est la résultante de la force de Lorentz magnétique qui s’exerce sur un volume $\d V = S\d\ell$ transverse au mouvement des particules, en fonction de l’intensité du courant électrique $I$ ? En déduire l’expression de la force de Laplace.
  
  \end{quest}
\end{exemple}
\examplesol{
  \begin{quest}
    \item	En appliquant cette expression à une charge $\d q = \rho \d V$ comprise dans un volume élémentaire $\d V$, on a
    $
    \Vec{\d F} = \d q\left(\vv{E}+\vv{v}\pv\vv{B}\right) = \left( \rho \Vec{E} + \rho \Vec{v} \pv \Vec{B}\right)\d V =  \left( \rho \Vec{E} + \Vec{\jmath} \pv \Vec{B}\right)\d V,
    $
    où $\Vec{\jmath}$ est la densité de courant électrique. On peut donc définir une force de Lorentz volumique
    \eqbox{
      \Vec{f} = \rho \Vec{E} + \Vec{\jmath} \pv \Vec{B}.
    }
  \end{quest}
}
  


\parag{Conservation de la charge}
%Il s'agit d'une loi capitale, comme toute équation de conservation elle nous dit quelque chose d'essentiel: la charge ne peut pas être créé ni détruite. Nous allons voir en fait que les équations de Maxwell contiennent implicitement l'équation de conservation de la charge.
%\subsubsection{Équation de conservation}
En régime variable, la densité volumique de charges $\rho$ (homogène à des \SI{}{C.m^{-3}}) et la densité de courant $\vv{j}$ (homogène à des \SI{}{A.m^{-2}}) ne sont  plus indépendants. De même que pour la masse $m$ en mécanique des fluides, nous souhaitons construire un théorie de l’électromagnétisme où la charge $q$ est conservée. La démonstration de l’équation de conservation est similaire.

\begin{methode}[Établir l’équation de conservation de la charge]
  \item	Définir comme système un volume élémentaire $\d V$ qui respecte la symétrie du problème :
  \begin{liste}
    \item tronçon de longueur $\d x$ en géométrie unidimensionnelle ; 
    \item coquille cylindrique de rayon $r$ et d’épaisseur $\d r$ en géométrie cylindrique ;
    \item coquille sphérique de rayon $r$ et d’épaisseur $\d r$ en géométrie sphérique. 
  \end{liste}
  \item Exprimer la variation de charge $\d q$ pendant la durée $\d t$ en fonction de $\dr{\rho}{t}$.
  \item Faire le bilan des charges entrantes et sortantes, en fonction des dérivées de $\j$.
  \item Simplifier et obtenir l’équation de conservation de la charge.
\end{methode}

\begin{enumerate}
  \item 
  Réalisons en effet un bilan de charges sur un tronçon de câble de section $S$ et de longueur $\d x$, donc de volume $S\d x$.
  
\item La variation de charge dans ce système ouvert pendant la durée $\d t$ est égal à la charge à l'instant $t+\d t$ moins la charge à l'instant $t$, soit
\begin{equation}\label{eq:1}
\d q(x,t) =\rho(x,t+\d t)S\d x-\rho(x,t)S\d x=\frac{\partial \rho}{\partial t}\d t S \d x.
\end{equation}
\item 
Cette variation de charge provient de ce qui est entré en $x$, moins ce qui est sorti en $x+\d x$, d'où
\begin{equation}\label{eq:2}
\d q(x,t)=\delta q_x(t)-\delta q_{x+\d x}(t)=j_x(x,t)S\d t-j_{x}(x+\d x,t)S\d t=-\frac{\partial j_x}{\partial x}\d x S \d t
\end{equation}
\item 
En égalisant les \Eq{1} et \Eq{2}, on obtient l'équation de conservation de la charge à $1$D,
% \begin{equation}
  \eqbox[]{
    \frac{\partial \rho}{\partial t}+\frac{\partial j_x}{\partial x}=0.
  }
% \end{equation}
\end{enumerate}
Elle se généralise à trois dimensions avec l’opérateur divergence.
\begin{theorem}[Équation de conservation de la charge]\label{form:q_cons}
\begin{equation*}
\frac{\partial \rho}{\partial t}+\dive{\j}=0
\end{equation*}
\end{theorem}
\noindent
Les postulats de l’électromagnétisme doivent être compatible avec l’équation de conservation de la charge.

\begin{exemple}{Bilan de charge dans un volume macroscopique}
On peut aussi démontrer la conservation de la charge avec une formulation intégrale. Soit un volume fixe $V$ contenant à l'instant $t$ la charge $Q(t)$.  Soit $\rho(M,t)$ la densité volumique de charge et $\j(M,t)$ la densité volumique de courant.
\begin{quest}
  \item	En découpant le volume $V$ en briques élémentaires de charge $\d q(t)$, montrer que 
  \begin{equation}
    \frac{\d Q}{\d t}=\iiint_{V}\frac{\partial\rho}{\partial t}\d V.
  \end{equation}
  \item Justifier que cette variation de charge s'écrit aussi
  \begin{equation}
    \frac{\d Q}{\d t}=-\oiint_{S}\j\cdot\vv{\d S}.
  \end{equation}
  \item En s'aidant du théorème de Green-Ostrogradski, démontrer la conservation de la charge.
\end{quest}

\end{exemple}
\examplesol{

On découpe par la pensée ce volume en briques élémentaires $\d V$ contenant la charge élémentaire $\d q(t)$ à l'instant $t$. Alors la variation de la charge contenue dans le volume $V$ par unité de temps s'écrit
\begin{equation}\label{eq:Q1}
\frac{\d Q(t)}{\d t}=\frac{\d}{\d t}\iiint_{V}\rho\,\d V=\iiint_{V}\frac{\partial\rho}{\partial t}\d V.
\end{equation}
Dans le même temps cette variation de charge par unité de temps dans le volume $V$ doit être égale à la charge \textit{reçue} par $V$ à travers la surface fermée $S$ qui délimite le contour de $V$, c'est-à-dire au flux de charge (donc l'intensité du courant) à travers cette surface compté dans le sens reçu
\begin{equation}\label{eq:Q2}
 \frac{\d Q(t)}{\d t}=-\oiint_{S}\j\cdot\vv{\d S} =  -\iiint_{V}\dive{\j}\d V.
 \end{equation}
 où le signe $-$ est important puisque $\vv{\d S}$ est orienté traditionnellement selon la normale extérieure à la surface. On a utilisé le théorème de Green-Ostrogradski. En égalisant les \Eq{Q1} et \Eq{Q2}, on en déduit que
 \begin{equation}
 \iiint_{V}\left(\frac{\partial \rho}{\partial t}+\dive{\j}\right)\d V=0
 \end{equation}
 pour tout volume $V$, d’où la \Formu{q_cons}.

}
%
%Mais cette équation est-elle compatible avec les équations de Maxwell ?

\subsection{Équations de Maxwell}
De même qu'on ne \guill{démontre} pas le principe fondamental de la dynamique en mécanique, on peut pas démontrer les équations de Maxwell. Il s'agit du \imp{postulat} de base de l'électromagnétisme, qui permet d'expliquer la totalité des phénomènes électromagnétiques connus.

\parag{Divergences des champs} Nous avons vu les équations de Maxwell-Gausse en électrostatique et de Maxwell-Thompson en magnétostatique. Ces deux équations restent valables \emph{en régime variable}.

\begin{theorem}[Équations de Maxwell avec divergence]\label{form:}
\begin{equation*}
\dive{\vv{E}}=\frac{\rho}{\varepsilon_0}~(\text{Maxwell-Gauss})\qquad
\dive{\vv{B}}=0~(\text{Maxwell-Thomson})
\end{equation*}
\end{theorem}
L’équation de Maxwell-Gauss impose une relation avec les \imp{sources} : le champ électrique est engendré par les charges électrique. L’équation de Maxwell-Thompson impose une relation de \imp{structure} : le champ magnétique est de divergence nulle.

\parag{Rotationnels des champs}
Nous avons vu que le champ électrostatique vérifie $\Rot{E} = \Vec{0}$, et le champ magnétostatique vérifie $\Rot{B} = \mu_0 \j$. Ces deux équations ne sont \imp{pas} valables en régime variable, et doivent être remplacées par les équations suivantes.
\begin{theorem}[Équations de Maxwell avec rotationnel]\label{form:}
\begin{equation*}
\rot{\vv{E}}=-\frac{\partial \vv{B}}{\partial t}~(\text{Maxwell-Faraday})\qquad
\rot{\vv{B}}=\mu_0\j+\mu_0\varepsilon_0\frac{\partial \vv{E}}{\partial t}~(\text{Maxwell-Ampère})
\end{equation*}
\end{theorem}
\noindent
L’équation de Maxwell-Faraday impose une relation de \imp{structure} du champ électromagnétique : la variation temporelle d’un champ magnétique créé un champ électrique. L’équation de Maxwell-Ampère impose une relation avec les \imp{sources} et une relation de \imp{structure} : le champ magnétique est créé par un courant électrique ou une variation temporelle du champ électrique.  En régime \imp{variable}, les équations aux rotationnels (Maxwell-Faraday et Maxwell-Ampère) couplent donc les champs $\vv{E}$ et $\vv{B}$ : c’est pourquoi il est nécessaire d’étudier le champ \imp{électromagnétique} dans sa globalité.

% Quelques remarques sur ces équations:
% \begin{liste}
% \item 

% \item

% \item En régime \imp{stationnaire} ($\frac{\partial }{\partial t}=0$), le champ électromagnétique se découple en une composante électrostatique créé par les charges ($\dive{\vv{E}}=\frac{\rho}{\varepsilon_0}$ et $\rot{\vv{E}}=0$), et une composante magnétostatique engendrée par les courants ($\dive{\vv{B}}=0$ et $\rot{\vv{B}}=\mu_0\vv{j}$).
%\item[$\bullet$] Le terme $\varepsilon_0 \frac{\partial \vv{E}}{\partial t}$ est homogène à une densité volumique de courant comme $\vv{j}$ (en \SI{}{A.m^{-2}}). On le note généralement $\vv{j_{\rm{D}}}$ et on l'appelle le courant de déplacement. Le nom est assez mal choisi, car il ne correspond physiquement pas à un déplacement de charges ! Il n'est donc pas physiquement analogue au courant de conduction $\vv{j}$, qui lui correspond véritablement à un déplacement des charges dans la matière ... Nous reviendrons sur ce point dans la suite en Sec.~\ref{subsubsec:ampère_généralisé} et dans le Complément~\ref{comp:condo}.
% \end{liste}

\subsection{Compatibilité avec la conservation de la charge}
\label{subsec:compatibilité_cons_charge}
On pourrait s’étonner de la présence du terme $\mu_0\varepsilon_0\frac{\partial \vv{E}}{\partial t}$ dans l’équation de  Maxwell-Ampère. 


\begin{remarque}
Historiquement, l'apport de Maxwell a été d’ajouter ce terme $\mu_0\varepsilon_0\frac{\partial \vv{E}}{\partial t}$ au terme $\mu_0\j$ déjà connu en magnétostatique.
% dans l'équation de Maxwell-Ampère afin de vérifier l'équation de conservation de la charge.
Cet apport est capital: c'est grâce à ce terme supplémentaire que Maxwell a prédit l'existence des ondes électromagnétiques, et a compris que la lumière est une onde électromagnétique. L'optique n'est donc qu'un cas particulier de l'électromagnétisme.
  
\end{remarque}
   Ce terme est en fait nécessaire pour assurer la conservation de la charge.
%, qui est un autre postulat de l’électromagnétisme.
On remarque que le vecteur $\varepsilon_0\frac{\partial \vv{E}}{\partial t}$
est homogène à une densité volumique de courant, tout comme $\j$ (en \SI{}{A.m^{-2}}). Cette grandeur représente le \imp{courant de déplacement}.
\begin{definition}[Courant de déplacement]\label{def:}
Le vecteur densité de \imp{courant de déplacement} (en \SI{}{A.m^{-2}}) est défini par
$$
\jd = \varepsilon_0\frac{\partial \vv{E}}{\partial t}.
$$
\end{definition}

L’équation de Maxwell-Ampère peut alors se récrire
\begin{equation}
\rot{\vv{B}}=\mu_0\vv{j}+\mu_0\jd = \mu_0\j_{\rm tot}.
\end{equation}


\begin{exemple}{Équations de Maxwell et conservation de la charge}
  En prenant la divergence de l'équation de Maxwell-Ampère, montrer que les équations de Maxwell impliquent la conservation de la charge :
   $$
   \dive{\j} + \frac{\partial \rho}{\partial t} = 0.
   $$
\end{exemple}
\examplesol{
Partons de l'équation de Maxwell-Ampère et prenons-en la divergence. Comme la divergence d'un rotationnel est nulle,
\begin{equation}
\dive{(\rot{\vv{B}})}=0=\mu_0\,(\dive{\j}+\dive{\jd}) \qquad\Rightarrow\qquad\dive{\j}+\dive{\jd}=0.
\end{equation}
En l’absence de courant de déplacement ($\jd = \Vec{0}$), on aurait $\dive{\j} = \Vec{0}$, et donc la charge ne serait pas conservée en régime variable ! Avec le courant de déplacement, sachant que les opérateurs vectoriels commutent avec les dérivées temporelles, il vient, d’après Maxwell-Gauss,
\begin{equation}
 0 = \dive{\j}+\dive{\jd} = \dive{\j} +  \varepsilon_0\frac{\partial }{\partial t}\underbrace{\dive{\vv{E}}}_{=\rho/\varepsilon_0} = \dive{\j} + \frac{\partial \rho}{\partial t}.
\end{equation}
}
% Le courant de déplacement est donc nécessaire pour assurer la conservation de la charge. Cela appelle quelques remarques.
% \begin{liste}
% \item
 L'équation de conservation de la charge est intrinsèquement incluse dans les équations de Maxwell ! En postulant les équations de Maxwell,  la charge électrique est automatiquement conservée.
% \item 
%, une conséquence des équations de Maxwell.
% \end{liste}
%C
%En régime dépendant du temps, le champ électromagnétique obéit aux équations de Maxwell suivantes (1856-1864)
%\begin{equation}
%\begin{cases}
%\dive{\vv{E}}=\frac{\rho}{\varepsilon_0}\qquad(\text{Maxwell-Gauss}),\\[2pt]
%\dive{\vv{B}}=0\qquad(\text{Maxwell-Thomson}),\\[2pt]
%\rot{\vv{E}}=-\frac{\partial \vv{B}}{\partial t}\qquad(\text{Maxwell-Faraday}),\\[2pt]
%\rot{\vv{B}}=\mu_0\vv{j}+\mu_0\varepsilon_0\frac{\partial \vv{E}}{\partial t}\qquad(\text{Maxwell-Ampère}).
%\end{cases}
%\end{equation}
%Quelques remarques sur ces équations:
%\begin{itemize}
%\item[$\bullet$] De même qu'on ne cherche pas à \guill{démontrer} le principe fondamental de la dynamique en mécanique, ici on peut pas démontrer les équations de Maxwell. Il s'agit du postulat de base de l'électromagnétisme, qui permet à ce jour d'expliquer la totalité des phénomènes électromagnétiques connus.
%\item[$\bullet$] En régime stationnaire ($\frac{\partial }{\partial t}=0$), remarquons que le champ électromagnétique se découple en une composante électrostatique créé par les charges ($\dive{\vv{E}}=\frac{\rho}{\varepsilon_0}$ et $\rot{\vv{E}}=0\Leftrightarrow\vv{E}=-\grad{V}$), et une composante magnétostatique engendrée par les courants ($\dive{\vv{B}}=0$ et $\rot{\vv{B}}=\mu_0\vv{j}$).
%\item[$\bullet$] En revanche en régime variable, les équations de Maxwell-Faraday et de Maxwell-Ampère couplent les champs $\vv{E}$ et $\vv{B}$, où une variation temporelle de champ magnétique engendre un champ électrique (équation de Maxwell-Faraday) et de même une variation temporelle de champ électrique engendre un champ magnétique (équation de Maxwell-Ampère). On parle donc du champ électromagnétique.
%\item[$\bullet$] Le terme $\varepsilon_0 \frac{\partial \vv{E}}{\partial t}$ est homogène à une densité volumique de courant comme $\vv{j}$ (en \SI{}{A.m^{-2}}). On le note généralement $\vv{j_{\rm{D}}}$ et on l'appelle le courant de déplacement. Le nom est assez mal choisi, car il ne correspond physiquement pas à un déplacement de charges ! Il n'est donc pas physiquement analogue au courant de conduction $\vv{j}$, qui lui correspond véritablement à un déplacement des charges dans la matière ... Nous reviendrons sur ce point dans la suite en Sec.~\ref{subsubsec:ampère_généralisé} et dans le Complément~\ref{comp:condo}.
%\end{itemize}


\section{Formulation intégrale de l’électromagnétisme}

Deux lois sont inchangée par rapport au cas statique : ce sont celles qui font intervenir le flux d'un champ à travers une surface fermée (Gauss et Thomson). Les autres lois sont modifiées par rapport à la statique : ce sont celles qui font intervenir la circulation d'un champ.

\subsection{Théorème de Gauss}
L'équation de Maxwell-Gauss étant inchangée en régime dépendant du temps, le théorème de Gauss reste valable. Notons $q_{\rm{int}}(t)=\iiint_{V}\rho(M,t)\d V$ la chargée emprisonnée à l'intérieur d’une surface fermée $S$ délimitant un volume $V$ à l'instant $t$.
\begin{theorem}[Théorème de Gauss]\label{form:gauss}
\begin{equation*}
\oiint_{M\in\Surf}\vv{E}(M,t)\cdot\vv{\d S}=\frac{q_{\rm{int}}(t)}{\varepsilon_0}
\end{equation*}
\end{theorem}

\begin{remarque}
  Ce théorème n'est pas aussi utile qu'en régime stationnaire, car l'étude des \imp{symétries est plus difficile} à mener en régime variable ! En effet, les symétries du champ $\vv{E}$ est  toujours dictée par les symétries de la densité volumique de charge $\rho$, mais aussi par celles du champ magnétique variable~$\vv{B}(t)$  d'après l'équation de Maxwell-Faraday.
Sauf dans le cas particulier de  l'ARQS électrique, qui ne figure pas au programme, les symétries et invariances des sources de $\vv{E}$ ne sont pas suffisantes pour choisir une surface de Gauss adapté. Vous serez donc guidé si vous devez utiliser le théorème de Gauss en régime variable.
% En pratique, vous n’aurez donc pas à utiliser le théorème de Gauss ne régime variable.
\end{remarque}

%conduisant à un calcul explicite du flux $\oiint_{M\in S}\vv{E}(M,t)\cdot\vv{\d S}$, même à un instant $t$ donné.

\subsection{Conservation du flux du champ magnétique}
L’équation de Maxwell-Thompson $\dive{\vv{B}}=0$ reste valable en régime variable. Considérons un tube de champ délimitée par deux sections $\Surf_1$ et $\Surf_2$, constituant une surface fermée de volume $V$. D'après le théorème de Green-Ostrogradski, on a 
\begin{equation}
\oiint_{M\in \text{tube}}\vv{B}(M,t)\cdot\vv{\d S}= \iiint_{M\in V} \dive{\vv{B}}(M,t)\,\d V= 0,
\end{equation}
soit
\begin{equation}\iint_{M\in\Surf_1}\vv{B}(M,t)\cdot\vv{\d S}=\iint_{M\in\Surf_2}\vv{B}(M,t)\cdot\vv{\d S} = \cte \text{ le long d’un tube de courant}.
\end{equation}
\begin{prop}[Conservation du flux magnétique]\label{prop:}
Le champ magnétique $\vv{B}(M,t)$ est à flux conservatif : le flux magnétique $\Phi_B$ est conservé le long d’un tube de champ à un instant donné $t$.
\end{prop}


\subsection{Loi de Faraday de l’induction}
Pour rappel, l’\imp{induction} électromagnétique désigne l’apparition d’une tension ou d’un courant à travers un circuit soumis à un flux magnétique variable. Il existe deux types d’inductions :
\begin{liste}
\item
Lorsqu’un circuit \imp{fixe} est soumis à un champ magnétique \imp{variable}, on parle d’induction de \imp{Neumann}. L’induction de Neumann a de nombreuses applications, vues en première année: le transformateur, les plaques à induction (chauffage par courants de Foucault), les badges magnétiques, etc.
\item Lorsqu’un circuit \imp{mobile} se déplace dans un champ magnétique \imp{stationnaire}, on parle d’induction de \imp{Lorentz}. L’induction de Lorentz intervient dans le moteur à courant continu, le freinage inductif des camions par courants de Foucault, le haut-parleur électrodynamique, les alternateurs et dynamos de vélo, etc.
\end{liste}


%%%%%%%%%%%%%%%%%%%%
\begin{figure}[b!]
\centering
\includegraphics[width=\textwidth]{induct.pdf}
\caption{(a) et (b) Induction de Neumann (champ variable mais circuit fixe). (a) Lorsque $B$ et donc $\Phi_{B}$ augmente, la fém $e$ est négative, donc $e$ et $i$ sont orientés dans le sens contraire à la règle de la main droite. La fém induite s'oppose à la variation de flux qui lui a donné naissance, selon la loi de Lenz. (c) et (d) Induction de Lorentz (champ magnétique constant et circuit mobile). \Femto }
\label{fig:induct}
\end{figure}
%%%%%%%%%%%%%%%%%

\parag{Démonstration de la loi de Faraday (induction de Neuman)}
Nous pouvons maintenant démontrer les lois de l’induction à partir de l’équation de Maxwell-Faraday. Considérons un circuit fermé fixe et indéformable $\Gamma$, entourant une surface $S$, soumis à un champ magnétique variable $\Vec{B}(M,t)$ (induction de Neuman). La loi de Maxwell-Faraday
\begin{equation}
\rot{\vv{E}} = -\dr{\Vec{B}}{t}
\end{equation}
indique qu’un champ électrique est créé dans l’espace, et en particulier le long du circuit. Le théorème de Stokes donne alors
%En régime variable, $\rot{\vv{E}}\neq 0$ et donc $\vv{E}\neq-\grad{V}$. Utilisons le théorème de Stokes avec l'équation de Maxwell-Faraday pour calculer la circulation de $\vv{E}$ sur un contour $\Gamma$ fixe, filiforme et fermé
\begin{align}
\oint_{M\in\Gamma}\vv{E}(M,t)\cdot\vv{\d l}& =\iint_{S}\rot{\vv{E}}\cdot\vv{\d S} \tag{Stokes}\\
&=-\iint_{S}\frac{\partial\vv{B}}{\partial t}\cdot\vv{\d S} \tag{Maxwell-Faraday} \\
&=-\frac{\d}{\d t}\iint_{S}\vv{B}\cdot\vv{\d S}. \tag{circuit fixe indéformable}
\end{align}
%où dans la dernière égalité nous avons utilisé que la surface $S(\Gamma)$ était indéformable pour pouvoir sortir la dérivée temporelle de l'intégrale.
On définit alors les deux quantités suivantes.
\begin{liste}
\item La \imp{force électromotrice} (fém, en $\SI{}{V}$)\footnote{Attention la force électromotrice n’est pas homogène à une force mais à une tension !} désigne la circulation de $\vv{E}$ sur un contour fermé
%En définissant
% \begin{equation}
  \eqbox[]{
    e=\oint_{\Gamma}\vv{E}\cdot\vv{\d l}.
  }
% \end{equation}
\begin{remarque}
  On voit bien que le champ électrique ne peut plus s'écrire comme un gradient du potentiel scalaire $V$ comme en régime variable, puisque sinon la circulation de $\E$ sur le circuit fermé $\Gamma$ serait nulle !
\end{remarque}
\item Le  \imp{flux magnétique} (en weber, $\SI{1}{W} = \SI{1}{T.m^2}$) désigne la grandeur
%à travers le circuit,
% \begin{equation}
  \eqbox[]{
    \Phi_B=\iint_{S}\vv{B}\cdot\vv{\d S}.
  }
% \end{equation}
\end{liste}
On en déduit alors la loi de Faraday de l'induction vue en première année. La loi de Faraday est également valable pour l’induction de Lorentz, mais nous n’en ferons pas la démonstration. 
\begin{theorem}[Loi de Faraday]\label{form:faraday}
  La variation du flux magnétique $\Phi_B$ à travers un circuit fermé induit une f.é.m.
\begin{equation*}
e=-\frac{\d\Phi_B}{\d t}
\end{equation*}
\end{theorem}
\noindent La \Formu{faraday} peut traduire deux effets.
\begin{liste}
\item Si $\Phi_B$ varie (soit $B$ varie : induction de Neuman, soit $S$ varie : induction de Lorentz), une fém $e$ est induite dans le circuit fermé. Si le milieu est conducteur, un courant électrique s’étblit. L’intensité du courant $i$ est orienté \important{dans le même sens} que $e$ (le circuit fonctionne comme un générateur).  Ces situations sont représentées sur la \Fig{induct}.
% apparaît dans le circuit\footnote{Il est \important{absolument indispensable} lors de l'étude d'un phénomène d'induction d'orienter le circuit, et donc de faire un schéma s'il n'est pas donné. Rappelons qu'il faut \important {toujours} calculer $\Phi_B$ avec le vecteur surface élémentaire orienté dans le sens du courant induit, c'est-à-dire dans le sens de la fém induite $e_{\rm{ind}}$.}. Pour changer ce flux, on peut soit faire varier le champ magnétique dans le temps et garder la surface fixe (Neumann), ou bien faire varier la surface dans le temps et garder le champ magnétique stationnaire (Lorentz).

\item Si une fém $e$ s’établit dans un circuit (un courant circule dans un conducteur), celle-ci induit un flux et donc un champ magnétique $B$.
% Ce champ magnétique induit va toujours s'opposer à la cause qui lui a donné naissance (loi de Lenz, le signe - dans la loi de Faraday).
\end{liste}
Dans les deux cas, la grandeur induite (fém $e$ ou champ $B$) s’\imp{oppose} à la grandeur qui lui a donné naissance (flux magnétique $\Phi_B$ ou courant d’intensité $i$), en raison du signe moins de la \Formu{faraday}.
%  Il s’agit de la loi de Lenz.
\begin{prop}[Loi de Lenz]\label{prop:}
La grandeur induite (fém $e$, champ $B$) s’oppose à celle qui lui donne naissance (flux $\Phi_B$, courant $i$).
\end{prop}
%
%Nous avons démontré cette relation en nous plaçant dans un référentiel particulier: celui où le circuit d'étude est fixe. On parle d'induction de Neumann: le circuit est fixe et le champ $\vv{B}(t)$ variable.
%
%
%Mais la loi de Faraday s'applique aussi si le champ magnétique ne dépend pas du temps mais que le circuit est mobile ! On parle d'induction de Lorentz. Ici aussi les applications sont nombreuses: le moteur à courant continu, le freinage inductif des camions par courants de Foucault, le haut-parleur électrodynamique, les alternateurs et dynamos de vélo ...
%
%Dans ce cas la surface $S(t)$ qui s'appuie sur $\Gamma$ dépend explicitement du temps, et il faut reprendre la démonstration précédente, mais le résultat final s'écrit toujours sous la forme $e=-\frac{\d\Phi_B}{\d t}$, c'est simplement que les expressions de $e$ et de $\Phi_B$ sont différentes selon le cas considéré
%\begin{equation}
%\begin{cases}
%e=\oint_{\Gamma}\vv{E}\cdot\vv{\d l}\quad\text{et}\quad\Phi_{B}=\iint_{S}\vv{B}(\JRouge{t})\cdot\vv{\d S} \quad\mbox{(Neumann)}, \\
%e=\oint_{\JRouge{\Gamma(t)}}(\vv{v_{\rm{e}}}\pv\vv{B})\cdot\vv{\d l}\quad\text{et}\quad\Phi_{B}=\iint_{\JRouge{S(t)}}\vv{B}\cdot\vv{\d S} \quad\mbox{(Lorentz)}.
%\end{cases}
%\end{equation}
%où nous avons noté $\vv{v_{\rm{e}}}$ la vitesse d'entraînement du circuit dans le référentiel d'étude (celui dans lequel $\vv{B}$ est stationnaire).
%
%\begin{figure}[h!]
%\centering
%\includegraphics[width=0.85\textwidth]{induction.pdf}
%\caption{Exemple simple d'induction de Neumann. Le circuit est fixe mais on fait varier le champ magnétique dans le temps de sorte que son flux à travers le circuit dépend du temps. Une fém induite apparaît qui s'oppose toujours à la variation de flux qui lui a donné naissance. Le champ magnétique induit est représenté en rouge, la fém induite et le courant induit en vert. En bas de la figure, le modèle électrocinétique équivalent.}
%\label{fig:induction}
%\end{figure}
%
%
%
%Dans le modèle électrocinétique équivalent, on prend en compte cette fém induite en rajoutant un générateur, avec un courant induit dans le même sens que la fém (Fig.~\ref{fig:induction}). Pour plus de détails sur l'induction électromagnétique, voir le cours de première année.

\begin{exemple}{Principe de l’alternateur}
Un alternateur est un convertisseur d’énergie mécanique en énergie électrique sous forme de courant alternatif. On en trouve, dans les centrales électriques, les éoliennes, les voitures, etc.
%%%%%%%%%%%%
  \begin{fig}[label]
    \centering
    %\capt{}
    \centering
      \begin{subfig}[label][0.45]
        \includegraphics[scale=1.5]{alternateur.pdf}
      \end{subfig}
      \begin{subfig}[label][0.45]
        \includegraphics[scale=0.7]{circuit.pdf}
      \end{subfig}
      \capt{(a) Modèle simplifié d’un alternateur \Femto~(b) Circuit équivalent. Le circuit se comporte comme un générateur de tension de fém $e$.}
    \end{fig}
%%%%%%%%%%%%%

On modélise un alternateur à induit mobile (rotor) par un cadre constitué de
$N$ enroulements d'un fil conducteur mis en rotation dans un champ magnétique uniforme et permanent. On suppose que le cadre est entraîné à une vitesse angulaire constante
$\omega$.
\begin{quest}
  \item	De quel type d’induction s’agit-il ? 
  \item Calculer le flux magnétique $\Phi_B$ à travers le circuit.
  \item En déduire la fém induite $e$.
\end{quest}

\end{exemple}
\examplesol{
\begin{quest}
  \item	
  Il s’agit de l’induction de Lorentz. 
  
  \item Le flux magnétique qui traverse le rotor vaut
  
  % $$
  \eqbox{
    \Phi_{B} = \iint_{S(t)}\vv{B}\cdot\vv{\d S} = NS \Vec{B} \cdot \Vec{n}(t) = NSB \cos(\omega t).
    }
% $$
\item 
La fém d’induction vaut, d’après la loi de Faraday,
% $$
\eqbox{
  e = -\dd{\Phi_B}{t} = NSB\omega \sin(\omega t).
  }
% $$
L'induit se comporte comme une source de tension alternative de pulsation $\omega$ et d'amplitude
$E
=
N
S
B \omega$.
\end{quest}

}


\subsection{Théorème d'Ampère généralisé}
\label{subsubsec:ampère_généralisé}
L'équation de Maxwell-Ampère est modifiée en régime variable par la présence du courant de le courant de déplacement. Essayons toutefois de trouver un théorème équivalent au théorème d'Ampère vu en statique, en utilisant le théorème de Stokes. Soit un contour fermé orienté $\Gamma$ fixe dans le référentiel d'étude, et $S$ une surface qui s'appuie sur $\Gamma$.

\begin{equation}
\oint_{\Gamma}\vv{B}(M,t)\cdot\vv{\d l}=\iint_{S}\rot{\vv{B}}\cdot\vv{\d S}=\mu_0\iint_{S}\left(\vv{j}+\varepsilon_0\frac{\partial \vv{E}}{\partial t}\right)\cdot\vv{\d S}.
\end{equation}
En séparant les deux termes et en intervertissant intégrale spatiale et dérivée temporelle (la surface $S$ est fixe), on obtient
\begin{equation}
  \oint_{\Gamma}\vv{B}(M,t)\cdot\vv{\d l} = \mu_0\iint_{S}\vv{j}\cdot\vv{\d S}+\mu_0\varepsilon_0\,\frac{\d}{\d t}\iint_{S}\vv{E}\cdot\vv{\d S} =  \mu_0 I_{\rm c} +\mu_0\varepsilon_0\,\frac{\d}{\d t}\iint_{S}\vv{E}\cdot\vv{\d S},
\end{equation}
où $I_{\rm c}$ désigne le courant de conduction enlacé.
En définissant le \imp{flux du champ électrique} à travers la surface
\begin{equation}
\Phi_E=\iint_{S}\vv{E}\cdot\vv{\d S},
\end{equation}
% le second terme s'écrit aussi
% \begin{equation}
%   \mu_0\varepsilon_0\,\frac{\d}{\d t}\iint_{S}\vv{E}\cdot\vv{\d S} = \mu_0\varepsilon_0\,\dd{\Phi_E}{t} = \mu_O
% \end{equation}
on obtient le théorème d'Ampère généralisé. On remarque que la grandeur
\begin{equation}
\varepsilon_0\frac{\d \Phi_{E}}{\d t} =  \iint_{S}\varepsilon_0\dr{\vv{E}}{t}\cdot\vv{\d S} = \iint_S \jd \ps \dS = I_{d}
\end{equation}
s'interprète comme l'intensité du courant de déplacement enlacé par le contour $\Gamma$.
\begin{theorem}[Théorème d'Ampère généralisé]\label{form:}
  La circulation de $\B$ sur un contour fermé est égal au courant total enlacé (conduction et déplacement)
$$
  \oint_{\Gamma}\vv{B}\cdot\vv{\d l}=\mu_0 I_{\Gamma} = \mu_0\left(I_{\rm{c}}+I_{\rm d}\right).
$$
\end{theorem}
\noindent

\begin{remarque}
  Attention, le courant de déplacement ne correspond pas à un déplacement réel de charges comme le courant de conduction !
\end{remarque}
% Insistons encore une fois,  mais à une variation temporelle du flux de $\vv{E}$ à travers la surface $S(\Gamma)$ qui s'appuie sur le contour d'Ampère $\Gamma$.
Lors de l'analyse des symétries pour le choix du contour d'Ampère, il est nécessaire d'analyser les symétries de $\vv{j}$ mais aussi de $\jd$. Dans l'immense majorité des cas, vous serez amenés à utiliser le théorème d'Ampère dans des régimes où le courant de déplacement est négligeable devant le courant de conduction.

\begin{remarque}
Nous verrons dans le chapitre suivant que dans un métal, le courant de conduction $\vv{j}$ et le courant de déplacement ne sont pas indépendants mais reliés par la loi d'Ohm locale $\vv{j}=\sigma\vv{E}$ avec $\sigma$ la conductivité électrique et donc en dérivant par rapport au temps $\frac{\partial \vv{j}}{\partial t}=\frac{\sigma}{\varepsilon_0}\vv{j_{\rm{D}}}$. Il en résulte que dans un métal et pour des fréquences d'étude raisonnables le courant de déplacement est négligeable devant le courant de conduction (voir le chapitre suivant).
\end{remarque}


\begin{exemple}{Le condensateur et la conservation du courant}
  On considère un condensateur plan idéal dont les armatures sont de surface $S$. En régime variable, on note $\E(z,t)$ le champ électrique entre les armatures et $Q(t)$ la charge sur l'armature positive.
  %Initialement le condensateur est chargé, on le place dans un circuit avec une résistance et on ferme l'interrupteur. On néglige les effets de bord pour le calcul du champ entre les armatures.

  \begin{figure}[H]
    \centering
    \includegraphics[width=0.49\textwidth]{capacitor.pdf}
    \caption{Application du théorème d'Ampère sur deux surfaces différentes.}
    \label{fig:capa}
    \end{figure}

\begin{quest}
  \item	
  Quelle est l'expression du champ $\E(z,t)$ entre les armatures du condensateur ? Justifier.
\end{quest}
Calculons la circulation de $\vv{B}$ le long du contour $\Gamma$ de la \Fig{capa}. On peut choisir pour cela deux types de surfaces : une surface $S_1$ qui coupe le fil; une surface surface $S_2$ qui passe entre les armatures sans couper le fil.
\begin{quest}
  \item Calculer la circulation de $\vv{B}$ le long du contour $\Gamma$ en utilisant le théorème d’Ampère et la surface $S_1$ pour calculer le courant total enlacé. 
  \item Calculer la circulation de $\vv{B}$ le long du contour $\Gamma$ en utilisant le théorème d’Ampère et la surface $S_2$ pour calculer le courant total enlacé. Commenter.
\end{quest}




  
\end{exemple}
\examplesol{
  \begin{quest}
    \item	L'application du théorème de Gauss généralisé donne la même expression qu'en statique, i.e.
    \begin{equation}
      \E(t)=\frac{Q(t)}{S\epz}\ez.
    \end{equation}
    \item Pour la surface $S_1$, le courant de déplacement est nul et le courant total enlacé est le courant de conduction $I_{\rm{e}}=I(t)$. On a donc 
    $$
    \oint_{\Gamma}\vv{B}\cdot\vv{\d l}=\mu_0 I(t).
    $$
    \item 
    Pour la surface $S_2$, le courant de conduction est nul et le courant total enlacé est la somme du courant  et du courant de déplacement $I_{\rm{enla}} = \varepsilon_0\frac{\d \Phi_E}{\d t}$, avec
    $$
    \Phi_E = \iint_{S_2}\vv{E}\cdot\vv{\d S} = \frac{Q(t)}{\varepsilon_0}.
    $$
    On a donc
    $$
    \oint_{\Gamma}\vv{B}\cdot\vv{\d l}=\mu_0\varepsilon_0\frac{\d \Phi_E}{\d t} = \mu_0 \dd{Q}{t}.
    $$
Les deux expressions sont bien égales en vertu de la conservation de la charge.
  \end{quest}
  
}

\subsection{Bilan}
\begin{figure}[h!]
\centering
\includegraphics[width=0.75\textwidth]{schema_bilan_maxwell.pdf}
\caption{Bilan des équations de Maxwell : lien entre les sources $(\rho,\j)$ et les champs ($\E,\B)$.}
\label{ref:tableau_bilan}
\end{figure}

Un résumé schématique du lien entre les sources du champ (charges et courants) et le champ électromagnétique créé par ces sources est représenté en Fig.~\ref{ref:tableau_bilan}.

\section{Énergie électromagnétique}
\subsection{Puissance cédée par le champ à la matière}
Le champ électromagnétique exerce sur une seule charge $q$ se déplaçant à la vitesse $\vv{v}$ la force de Lorentz $\vv{F}=q(\vv{E}+\vv{v}\pv\vv{B})$.
%Puisque la composante magnétique est orthogonale à la vitesse, elle ne contribue pas à la puissance de la force. En effet, un champ magnétique ne peut pas accélérer les charges contrairement au champ électrique, mais il peut dévier leur trajectoire.
La puissance cédée par le champ électromagnétique à une particule de charge $q$ s'écrit donc
\begin{equation}
\mc{P}=q(\vv{E}+\vv{v}\pv\vv{B})\ps \Vec{v} = q\vv{E}\cdot\vv{v}.
\end{equation}
Soit un milieu contenant une densité de particules chargées $n$ identiques de charge $q$. La puissance cédée un volume $\d V$ de matière vaut
\begin{equation}
  \d \mc{P} =\mc{P} \times (n\d V) = nq\vv{v} \ps \vv{E} \d V = \Vec{j} \ps \Vec{E} \d V.
\end{equation}
On voit apparaître le vecteur densité de courant de charge $\j = \rho \Vec{v} = n q \Vec{v}$. Ce résultat se généralise à un milieu contenant différents porteurs de charges $i$, de densité volumique de particules $n_i$, de charge $q_i$ et se déplaçant à la vitesse $\vv{v_i}$, avec dans ce cas  $\vv{j}=\sum_{i}q_i n_{i}\vv{v_i}=\sum_{i}\rho_{i}\vv{v_i}$.

\begin{theorem}[Puissance volumique cédée à la matière]\label{form:}
La puissance volumique cédée par le champ à la matière
 s'écrit
 $$
 p_{\rm{vol}} = \j\cdot\vv{E}.
 $$
\end{theorem}
% En conclusion,et il suffit d'intégrer sur le volume $V$ pour obtenir la puissance totale cédée par le champ à la matière
% \begin{equation}
% \mc{P}=\iiint_{V}\vv{j}\cdot\vv{E}\,\d V.
% \end{equation}

\subsection{Équation locale de Poynting}
Si le champ électromagnétique est capable de céder de l'énergie à la matière, c'est qu'il en possède en premier lieu. Nous avons déjà vu l'expression des densités volumiques d'énergie électrique $u_e = \frac{1}{2}\epz E^2$ et magnétique $u_m = \frac{B^2}{2\mu_0}$. Cherchons à établir l'équation de conservation de l'énergie du champ électromagnétique.

\parag{Équation de conservation de l’énergie électromagnétique}

Comme point de départ, réexprimons la densité volumique de puissance $\vv{j}\cdot\vv{E}$ en fonction des champs uniquement. On remplace $\vv{j}$ avec l'équation de Maxwell-Ampère,
\begin{equation}
\j =\frac{1}{\mu_0}\rot{\vv{B}}-\varepsilon_0\frac{\partial \vv{E}}{\partial t}
\end{equation}
ainsi
\begin{equation}
\j \ps \E =\frac{1}{\mu_0}\vv{E}\cdot\rot{\vv{B}}-\frac{\partial}{\partial t}\left(\frac{1}{2}\varepsilon_0 E^{2}\right) =  \frac{1}{\mu_0}\vv{E}\cdot\rot{\vv{B}}-\frac{\partial u_e}{\partial t}.
\end{equation}
Pour faire apparaître l'énergie volumique magnétique, on utile la formule
\begin{equation}
  \dive{(\vv{E}\pv\vv{B})}=\vv{B}\cdot\rot{\vv{E}}-\vv{E}\cdot\rot{\vv{B}},
\end{equation}
qui combiné avec l'équation de Maxwell-Faraday $\rot{\vv{E}}=-\frac{\partial \vv{B}}{\partial t}$, donne
\begin{equation}
  \vv{E}\cdot\rot{\vv{B}} = -\frac{\partial}{\partial t}\left(\frac{B^2}{2}\right)-\dive{(\vv{E}\pv\vv{B})}  = -\mu_0 \dr{u_m}{t} - \dive{(\vv{E}\pv\vv{B})}  .
\end{equation}
On obtient donc
\begin{equation}
  \dive{\left(\frac{\vv{E}\pv\vv{B}}{\mu_0}\right)}+\dr{}{t}\left(u_e+u_m\right) = -\j\cdot\vv{E}.
\end{equation}
On définit alors
\begin{liste}
  \item La \imp{densité volumique d'énergie électromagnétique} (en $\SI{}{J.m^{-3}}$)
  % \begin{equation}
    \eqbox[]{
      u_{\rm em} = u_{\rm e}+u_{\rm m} = \frac{1}{2}\varepsilon_0 E^2+\frac{B^2}{2\mu_0}.
    }
  % \end{equation}
  \item Le \imp{vecteur de Poynting} (en $\SI{}{W.m^{-2}}$)
  % \begin{equation}
    \eqbox[]{
      \Poy = \frac{\vv{E}\pv\vv{B}}{\mu_0}.
      }
  % \end{equation}
  Il représente le vecteur densité de courant d'énergie électromagnétique.
\end{liste}

\begin{theorem}[Équation de convervation de l'énergie électromagnétique]\label{form:}
  $$
\dr{u_{em}}{t} + \dive \Poy = -p_{\rm vol}.
$$
\end{theorem}
L'énergie électromagnétique est donc transportée par le vecteur de  Poynting $\Poy$ et perdue au charge selon un terme de perte volumique $-\j \ps \E$.
% En réarrangeant les termes, on peut mettre cette équation sous la forme
% \begin{equation}
% \frac{\partial u_{\rm{em}}}{\partial t}+\dive{\,\vv{\Pi}}=-\mc{P}_{\rm{vol}}
% \end{equation}
% où nous avons défini
% \begin{itemize}
% \item[$\bullet$] $u_{\rm{em}}=\frac{1}{2}\varepsilon_0 E^2+\frac{1}{2\mu_0}B^2$ la densité volumique d'énergie du champ électromagnétique en $\SI{}{J.m^{-3}}$,
% \item[$\bullet$] $$ le vecteur de Poynting, qui correspond à un vecteur densité de courant d'énergie volumique, homogène à une énergie par unité de surface et par unité de temps, donc une puissance surfacique en
% \item[$\bullet$] $-\vv{j}\cdot\vv{E}$ la densité volumique de puissance reçue par le champ électromagnétique de la part de la distribution de courants de densité volumique $\vv{j}$.
% \item[$\bullet$] réciproquement $+\vv{j}\cdot\vv{E}$ désigne celle reçue par la matière. C'est elle qui dans un conducteur va être convertie sous forme de chaleur (effet Joule).
% \end{itemize}

\parag{Forme intégrale}
Pour interpréter cette équation il est possible de l'écrire sous forme intégrale. Soit un volume fixe $V$ délimité par une surface (fermée) S. D'après le théorème de Green-Ostrogradski,
\begin{equation}
\iiint_{V}\frac{\partial u_{\rm{em}}}{\partial t}\d V=-\iiint_{S}\dive \vv{\Pi}\d V- \iiint_{V}\vv{j}\cdot\vv{E}\,\d V = -\oiint_{S}\vv{\Pi}\cdot\vv{\d S} - \iiint_{V}\vv{j}\cdot\vv{E}\,\d V.
\end{equation}
On notant
\begin{equation}
  U_{\rm em} = \iiint_{V} u_{\rm{em}}\d V =  \iiint_{V} \left( \frac{1}{2}\varepsilon_0 E^2+\frac{B^2}{2\mu_0} \right)\d V
\end{equation}
l'énergie électromagnétique contenue dans le volume $V$, on obtient
% \begin{equation}
\eqbox[]{
  \dd{U_{\rm em}}{t} = -\oiint_{S}\vv{\Pi}\cdot\vv{\d S} - \iiint_{V}\vv{j}\cdot\vv{E}\,\d V.
  }
% \end{equation}
La grandeur $\oiint_{S}\vv{\Pi}\cdot\vv{\d S} $ représente la puissance transportée à travers la surface $S$ et  $\iiint_{V}\vv{j}\cdot\vv{E}\,\d V $ la puissance cédée à la matière dans le volume $V$.


\section{Approximation des régimes quasi-stationnaires (ARQS)}
\label{subsec:ARQS}

% En électrocinétique, nous étudions les circuits en terme de potentiel $V$ et de champ électrique $E$, pulsation

Nous étudions maintenant ce qu'il se passe dans l'approximation des régimes lentement variables dans le temps, c'est-à-dire quasi-stationnaires (ARQS). Les champs sont créés par des sources, les charges et les courants, qui dépendent en principe du temps. Que se passe-t-il si les variations dans le temps d'un type de source (par exemple les courants) sont grands devant l'autre (du coup les charges) ? Comment peut-on simplifier les équations de Maxwell ?
\begin{remarque}
En toute rigueur, il existe deux ARQS: l'ARQS électrique, où les effets temporels des charges dominent ceux des courants; et l'ARQS magnétique, où les effets temporels des courants dominent ceux des charges. Conformément au programme, nous ne discuterons que le dernier cas, appelé uniquement ARQS dans la suite.
\end{remarque}

\subsection{Cadre de l’approximation}
Nous verrons que dans les ondes électromagnétiques se propage à une vitesse de l’ordre de la célérité de la lumière $c \simeq \SI{3e8}{m.s^{-1}}$. Cela signifie que sur un circuit électrique de taille caractéristique $L$, le champ électromagnétique met un temps certain $\tau$ pour se propager sur toute la longueur du circuit, d’expression
\begin{equation}
  \tau = \frac{L}{c}
\end{equation}

% Or les générateurs usuels (GBF) imposent un champ électromagnétique de période au mieux $T \sim \SI{e-6}{s} \gg \frac{L}{c}$ (fréquence maximale de $f \sim \SI{1}{MHz}$) : on peut considérer que le champ électromagnétique se propage le long du circuit de façon quasi-instantannée.
\begin{exemple}{ARQS}
  Calculer le temps de propagation $\tau$ du champs électromagnétiques sur un circuit de taille $L$ usuelle. Comparer à la période $T$ d'un signal électrique de fréquence  usuelle pour les GBF de TP.
\end{exemple}
\examplesol{
  Pour un circuit de taille $L \sim \SI{10}{cm}$, le temps de propagation vaut 
  \eqbox{
  \tau  \sim \SI{3e-8}{s}.
  }
  La période d’un signal de fréquence $f \sim \SI{1}{MHz}$ vaut
\eqbox{
  T \sim \SI{1e-6}{s}.
}
On voit donc qu’en pratique, on a 
\eqbox{
  \tau \ll T.
}
La propagation de l’onde est quasi instantanée, et les champs électromagnétiques sont uniformes à l’échelle du circuit.
}
%Le champ peut donc être considéré come quasi-stationnaire ()

 % sur une distance la vitesse de l’information transportée par le signal électrique   L'approximation des régimes quasi-stationnaires consiste à supposer que l'onde se propage quasi-instantanément ($c\rightarrow \infty$), donc à négliger le retard dû à la propagation. Cette approximation est valable lorsqu'on étudie des phénomènes à des distances $d$ des sources telles que
\begin{prop}[ARQS]\label{prop:}
  L’Approximation des Régimes Quasi Stationnaires  est valable lorsque la taille du circuit vérifie
  \begin{equation*}
  L\ll cT.
  \end{equation*}
\end{prop}

% où $T$ désigne le temps caractéristique d'évolution des sources du champ électromagnétique. Par conséquent, les champs $\vv{E}$ et $\vv{B}$ sont uniformes à l'échelle du dispositif.

% \examplesol{
%   En régime sinusoïdal, $cT=\lambda$ la longueur d'onde de l'onde électromagnétique créé par les sources. Dans l’ARQS, la longueur caractéristique du dispositif est négligeable devant la longueur d'onde de l'onde émise
%   % \begin{equation}
%     \eqbox{
%       L\ll \lambda.
%     }
%     % \end{equation}
% }



% \begin{exemple}
% Typiquement lors de l'étude des phénomènes électriques en travaux pratiques, $\lambda=c/f<\SI{300}{m}$ pour $f=\SI{1}{MHz}$. La taille typique d'un circuit étant négligeable devant la longueur d'onde, on se trouve bien dans le régime de l'ARQS. C'est d'ailleurs pour cela qu'en électronique on utilise ce qu'on qualifie de générateur basse fréquence (GBF) : même à \SI{1}{MHz}, leur fréquence est négligeable devant $c/d$ ce qui nous place dans le cadre de l'ARQS.
% \end{exemple}

\subsection{Lois locales et intégrales dans l’ARQS}

\parag{Équations de Maxwell} Évaluons l’ordre de grandeur de différents termes des équations de Maxwell dans l’ARQS. Soit $L$ la distance typique (taille du circuit) et $T$ le temps typique  (période).

\begin{enumerate}
  \item \udl{Maxwell-Faraday} : les amplitudes des champs $E$ et $B$ vérifient
  \begin{equation}
    \rot\E = -\dr{\B}{t} \Rightarrow \frac{E}{L} \sim \frac{B}{T} \Rightarrow \boxed{B \sim \frac{E T}{L}}.
   \end{equation}
   \item \udl{Maxwell-Gauss} : la densité de charge volumique $\rho$ est liée au champ électrique par
    \begin{equation}
      \dive\E = \frac{\rho}{\epz} \Rightarrow \frac{E}{L} \sim \frac{\rho}{\epz} \Rightarrow \boxed{\rho \sim \frac{\epz E}{L}}.
    \end{equation}
    \item \udl{Maxwell-Thomson} : l’équation de Maxwell-Thomson est inchangée,  soit
    % \begin{equation}
      $\dive\B = 0.$
    % \end{equation}
   \item \udl{Maxwell-Ampère} : le terme issu du courant de déplacement vaut
   \begin{equation}
     \mu_0 \epz \dr{\E}{t} \sim \frac{E}{c^2 T}.
   \end{equation}
   On peut le comparer au terme magnétique,  
   \begin{equation}
     \frac{\norme{\mu_0 \epz \dr{\E}{t}}}{\norme{\rot\B}} \sim \frac{E/c^2T}{\frac{B}{L}} \sim \left(\frac{L}{cT}\right)^2 \ll 1
   \end{equation}
   dans l’ARQS. Le courant de déplacement est donc négligeable devant le courant de conduction, où de façon équivalente, le champ magnétique domine le champ électrique. On parle d’\imp{ARQS magnétique}. L’équation de Maxwell-Ampère  a donc même forme qu’en statique. 
\end{enumerate}

% L'équation de Maxwell-Ampère,
% \begin{equation}
% \rot{\vv{B}}=\mu_0\j+\mu_0\epz\frac{\partial \vv{E}}{\partial t}
% \end{equation}
% se simplifie donc dans l'ARQS par
% \begin{equation}
% \rot{\vv{B}}=\mu_0\j.
% \end{equation}
% Il s'agit en fait de l'équation de Maxwell-Ampère de la magnétostatique ! Le courant de conduction $\j$ domine le courant de déplacement $\j_{\rm d}$.
% Cela revient àoù nous avons négligé le courant de déplacement devant le courant de conduction.

\begin{theorem}[Équation de Maxwell-Ampère dans l’ARQS]\label{form:}
L’équation de \imp{Maxwell-Ampère} prend la même forme dans l’ARQS qu’en statique,
$$
\rot{\vv{B}}=\mu_0\j.
$$
% Les autres équations de Maxwell sont inchangées.
La forme intégrale de cette équation est le \imp{théorème d’Ampère}.

\end{theorem}

\parag{Lois de Kirchoff} Dans le cadre de l’ARQS, nous pouvons démontrer les lois de Kirchoff.

\begin{enumerate}
  \item \udl{Loi des nœuds}.
  
  La conservation de la charge s’obtient en prenant la divergence de l'équation de Maxwell-Ampère,
  % \begin{equation}
  \eqbox{
  \dive{\j}=0,
  }
  % \end{equation}
où le terme $\frac{\partial \rho}{\partial t}$ n'apparaît plus. 
L’effet des variations temporelles de charges est donc négligeable devant celui des courants. 
% L'équation $\dive{\j}=0$ est simplement l'expression locale de la loi des nœuds.
\begin{remarque}
  On peut retrouver cette équation remarquant qu’en ordre de grandeur,
\begin{equation}
\frac{\partial \rho}{\partial t} \sim \frac{\rho}{T} \sim \frac{\epz E}{L T} \quad \text{ et } \quad \dive{\j} \sim \frac{B}{\mu_0 L} \sim \frac{E}{\mu_0 L^2},
\end{equation}
et donc 
\begin{equation}
\frac{\partial{\rho} / {\partial t}}{\dive{\j}} \sim \left(\frac{L}{cT}\right)^2 \ll 1.
\end{equation}
\end{remarque}

\begin{theorem}[Loi des nœuds]\label{form:}
L’équation de \imp{conservation de la charge} dans l’ARQS prend la même forme qu’en statique,
$$
\dive \j = \Vec{0}.
$$
La forme intégrale de cette équation est la \imp{loi des nœuds}.
\end{theorem}

\item \udl{Loi des mailles}.

Considérons une maille, i.e. un circuit fermé orienté $\Gamma$, constituée de $N$ branches.
En prenant la circulation de l’équation de Maxwell-Faraday, on trouve la loi de Faraday de l’induction,
\begin{equation}
  e = \oint_{\Gamma}\E\cdot\vv{\d l} = -\frac{\d}{\d t}\iint_{S}\B\cdot\vv{\d S} = -\frac{\d \Phi_B}{\d t},
\end{equation}
Or la circulation de $\E$ est aussi égale à la somme algébrique des tensions $u_i$ dans le circuit, soit
\begin{equation}
  \oint_{\Gamma}\E\cdot\vv{\d l} = \sum_{i=1}^N u_i = e.
\end{equation}
On peut ajouter dans la maille un générateur fictif de tension $e = - \dd{\Phi_B}{t}$, ou un récepteur fictif de tension $u_0 = \dd{\Phi_B}{t}$, tension que l’on inclut dans la somme algébrique des tensions. On obtient alors la loi des mailles,
% $$
\begin{equation}
  \sum_{i=0}^N u_i = 0.
\end{equation}
% $$

\begin{prop}[Loi des mailles]\label{prop:}
  L’équation de \imp{Maxwell-Faraday} dans l’ARQS prend la même forme qu’en régime variable,
$$
\rot \E = - \dr{\B}{t}.
$$
La forme intégrale de cette équation est la \imp{loi des mailles}.
\end{prop}

\end{enumerate}

\parag{Caractéristiques des composants linéaires}

\begin{enumerate}
  \item \udl{Bobine à induction}.
  
  Considérons un solénoïde d’inductance propre $L$ dans l’ARQS. Le champ magnétique et le flux propre ont la même expression qu’en statique, et donc en particulier, 
  \begin{equation}
  \Phi = L i.
  \end{equation}
  La loi de Faraday de l’induction s’applique comme en régime variable dans l’ARQS, et donc
  \eqbox{
  u_L = -e = \dd{\Phi}{t} =  L \dd{i}{t}.
  }
\end{enumerate}

En conclusion, dans l’ARQS :
\begin{liste}
  \item Les champs électriques et magnétiques se calculent comme en \imp{statique} avec les théorèmes de Gauss et d'Ampère.
  \begin{remarque}
    Le théorème de Gauss est nécessaire pour décrire la caractéristique $i=C\dd{u}{t}$ d’un condensateur. L’équation de Maxwell-Ampère en statique est nécessaire pour que la loi des nœuds soit valable.
  \end{remarque}
  % \begin{remarque}
  %   Le théorème de Gauss, valable en régime quelconque, permet de décrire la caractéristique $i=C\dd{u}{t}$ d’un condensateur.
  % \end{remarque}
  \item La loi de Faraday de l’induction $e=-\frac{\d \Phi}{\d t}$ s'applique comme en régime \imp{variable}.
\begin{remarque}
    La loi de Faraday est nécessaire pour décrire la caractéristique $u=L\dd{i}{t}$ d’une bobine à induction.
\end{remarque}

\end{liste}

\begin{remarque}
  Comme le terme magnétique domine le terme électrique dans l’équation de Maxwell-Ampère, on parle d’ARQS \guill{magnétique}. Il existe un autre régime d’ARQS, l’ARQS \guill{électrique}, qui décrit notamment le fonctionnement d’un condensateur en régime variable.
\end{remarque}

\begin{exemple}{Lois de l’ARQS en électrocinétique.}
  \begin{quest}
    \item	Montrer que dans l’ARQS, le théorème de Gauss assure que la caractéristique d’un condensateur vérifie $i=C\dd{u}{t}$.
    \item Montrer que dans l’ARQS, la loi de Faraday assure que la caractéristique d’une bobine à induction vérifie $u=L\dd{i}{t}$.
    \item Montrer que dans l’ARQS, la conservation de la charge assure la loi des nœuds. Dans quelle mesure la loi des mailles est-elle valide ?
  \end{quest}
\end{exemple}
\examplesol{
  \begin{quest}
    \item	La loi de l’induction donne $e=-\frac{\d \Phi}{\d t}$. Or dans une bobine à induction dans l’ARQS, le champ magnétique se calcule comme en statique, et donc $\Phi=Li$ avec $L$ l’auto-inductance.   On en déduit que $e=-L\dd{i}{t}$, soit $u=L\dd{i}{t}$ en convention récepteur.
  \end{quest}
}
% Toutes les autres équations de Maxwell sont inchangées par rapport au régime quelconque. Pour l'équation de Maxwell-Gauss et de Maxwell-Ampère tout se passe comme en électrostatique et en magnétostatique respectivement, on peut donc utiliser les formes simples des théorèmes de Gauss et d'Ampère. De plus, la loi de Faraday $e=-\frac{\d \Phi}{\d t}$ s'applique toujours comme en régime variable, c'est le domaine d'étude de l'induction, comme décrite en première année.



% \newpage
% \begin{subappendices}
% \section{Le condensateur et la conservation du courant}
% \label{comp:condo}
% On considère un condensateur plan dont les armatures circulaires ont pour surface $S$. Initialement le condensateur est chargé, on le place dans un circuit avec une résistance et on ferme l'interrupteur. On néglige les effets de bord pour le calcul du champ entre les armatures.
%
% \begin{figure}[h!]
% \centering
% \includegraphics[width=0.7\textwidth]{condo_jD.pdf}
% \end{figure}
%
% \begin{enumerate}
% \item Donner l'expression du champ $\vv{E}(t)$ entre les plaques du condensateur.
% \item En considérant pour contour $\Gamma$ un petit cercle centré sur le fil, calculer en négligeant les effets de bord la circulation de $\vv{B}$ en choisissant une surface $S_1$ s'appuyant sur $\Gamma$ qui coupe le fil sans passer par l'intérieur du condensateur, puis une surface $S_2$ qui passe entre les armatures sans couper le fil. Commenter.
% \end{enumerate}
% \textit{Solution:}
% 1. $\vv{E}(t)=\frac{Q(t)}{S\varepsilon_0}\vv{e_z}$ par analogie\footnote{Peut-on analyser les invariances et symétries comme en statique puisque d'après l'équation de Maxwell-Faraday un champ magnétique dépendant du temps est aussi source de $\vv{E}$ ? La réponse est ici oui, car dans un condensateur $\rot{\vv{E}}=0$ même en régime variable. On peut donc utiliser les mêmes invariances et symétries et finalement le même théorème de Gauss. Ce résultat peut se démontrer dans le cadre de l'ARQS électrique, dont l'étude est hors programme. Seule l'ARQS magnétique figure au programme (Sec.~\ref{subsec:ARQS}).} avec le cas statique.
%
% 2. Entre les plaques du condensateur, $\vv{j}=0$ car il n'y a pas de charges mobiles (il y a du vide entre les plaques), donc pas de courant de conduction. Dans le fil qui relie les armatures c'est l'inverse: on néglige le courant de déplacement et il n'existe qu'un courant de conduction. Autrement dit le théorème d'Ampère s'écrit
% \begin{equation}
% \begin{split}
% \oint_{\Gamma}\vv{B}\cdot\vv{\d l}&=\mu_0\iint_{S(\Gamma)}\left(\vv{j}+\varepsilon_0\frac{\partial \vv{E}}{\partial t}\right)\cdot\vv{\d S}\\
%  &=\begin{cases}
%  \mu_0\iint_{S_1} \vv{j}\cdot\vv{\d S}=\mu_0 i(t) \\
%  \mu_0\iint_{S_2}  \varepsilon_0\frac{\partial \vv{E}}{\partial t}\cdot\vv{\d S}=\mu_0\varepsilon_0\frac{\d }{\d t}\left(\iint_{S_2}\frac{Q(t)}{\varepsilon_0 S}\vv{e_z}\cdot\vv{\d S}\right)=\mu_0\frac{\d Q}{\d t}.
%  \end{cases}
%  \end{split}
%  \end{equation}
% Le courant de déplacement est donc essentiel pour retrouver le fait que $i(t)=\frac{\d Q(t)}{\d t}$ à savoir que l'intensité du courant qui rentre et qui sort du condensateur est égale à la variation temporelle de la charge $Q(t)$ de ses armatures.
% Fig

% \end{subappendices}


\end{document}
